\section{Bridges to other areas, and the directions they open}\label{sec:bridges}

Much of the programme consists of bridges between the zeta function and other parts of mathematics, and in the author's view they are among its most valuable outcomes. This section treats each bridge in the same way. It first gives the typed map: what is mapped to what, and what crosses. It then says what is established, what the bridge suggests on each side, and what my assessment is. The assessment is a view held at this stage and is labelled as such. Each item says whether it is a theorem, an identity of objects, an analogy or a proposal.

\subsection{Meyer's quotient and interpolation theory}

The jet map $J:Q\to\prod_\rho A_\rho$ carries Meyer's quotient to the jets at the zeros, and what crosses is the interpolation problem at the zeros. It is established that $J$ is injective with dense image in the space $\Lambda_m$ of rapidly decreasing jet tuples (Corollary~\ref{cor:jetimage}), and that it is onto exactly when the principal parts of $1/\zeta$ at the zeros are polynomially bounded (Theorem~\ref{thm:primary}).
The convergence of the primary decomposition is thereby the same kind of question as Berenstein and Taylor's interpolating varieties \cite{BerensteinTaylor} and Bondarenko, Radchenko and Seip's Fourier interpolation at the zeta zeros \cite{BondarenkoRadchenkoSeip}. The Gonek--Hejhal conjectures concern the averages $\sum|\zeta'(\rho)|^{-2k}$, while Theorem~\ref{thm:primary} asks the pointwise question (an analogy of questions; the theorem itself is proved).

For interpolation theory the bridge offers the zeta zeros as a test case for interpolating varieties; for $\zeta$ it offers a structural reason to care about pointwise lower bounds for $|\zeta'(\rho)|$.
My assessment: this is the most precise translation in the programme, because it turns a structural property of its central object into a classical analytic condition. The case of multiple zeros, where in the weighted spaces $A_p$ the leading coefficients alone suffice \cite[Theorem~4]{BerensteinTaylor}, is the natural next question (Question~\ref{q:primary}).

\subsection{The two-line system, Weil's explicit formula and Turing's method}

Three maps make up this bridge. The decomposition $f\mapsto(u,v)$ into parts even and odd under the translation carries the chiral form to $Q_A(f)=2\Omega(u)-2\Omega(v)$. The loop $T^{-1}A$ is Weil's reflection $\#$ through $\tfrac12$. And the holonomy count goes to Turing's index $(N(T)-V(T))/2$.
All three are theorems (Section~\ref{sec:twolines}). The mirror through the pole at $0$ followed by the passage between the lines is Weil's reflection. On the diagonal of the even data the chiral form is twice Weil's form (Proposition~\ref{prop:chiral}). The count of nontrivial holonomies plus $\sum\lfloor m_\rho/2\rfloor$ over the on-line zeros is Turing's index (Proposition~\ref{prop:index}).
On Weil's side the bridge suggests a test family $\{f_k\}$ that is even under the translation by construction and has the explicit prime side $-G'/G$ (Corollary~\ref{cor:chiral}). On Turing's side it suggests that an index for off-line zeros alone needs a map that separates off-line pairs from multiple zeros on the line.

My assessment: the bridge recasts Weil's positivity as a signature statement about a form that is hyperbolic for every configuration of zeros. Since the whole form has the same signature with or without RH, I expect progress along this bridge to come from the diagonal family, where the positivity question is concentrated (Question~\ref{q:diagonal}).

\subsection{The programme's weights and Deligne's Weil~II}

This bridge carries four maps:
\begin{itemize}
\item a zero $\rho$ goes to the eigenvalue $p^\rho$, of weight $2\re\rho$;
\item the even tensor powers go to the positivity step of Weil~II~1.5;
\item the timed-prime measure goes to Deligne's positive measure;
\item the ramification operator $P_b$ goes to a relation $FNF^{-1}=Q^{-1}N$ of Deligne's form.
\end{itemize}
Established are the results of Section~\ref{sec:deligne}, and further that Weil~II~§2.1 specialised to $\zeta$ is the classical proof of $\zeta(1+it)\ne0$, with the timed-prime measure mapping onto Deligne's positive measure (\lbl{DB9}, \lbl{DR0--DR4}; theorems).
In the programme it is ramification, not counting, that satisfies a relation of Deligne's form ($P_bNP_b^{-1}=b^{-1}N$ on the image of $P_b$). Every solution of the relation on a jet block forces the weight shift $-2j$ of the programme's regrading, and such relations have solutions for every nilpotent, so what is established is a comparison of relations (\note{35\_} §3).
Under that regrading the jet block of a zero of multiplicity $m$ satisfies the conclusion of Weil~II~1.8.4 with centre $2\re\rho-(m-1)$. For $m\ge2$ the block is mixed in the sense of Weil~II (1.7.4), and it is pure in the sense later used for Weil--Deligne representations \cite{TaylorYoshida}. The regrading is chosen in the programme rather than derived from an arithmetic source, so this is a comparison of structures (\note{36\_} Remark 36.R1).
Deligne's step $\alpha\mapsto\alpha^2$ ($H^1\otimes H^1\hookrightarrow H^2$) has as additive shadow Kurokawa's tensor product of zeta functions, with spectral parameters $\rho+\rho'$ (Manin \cite{Manin}, §1.4). Weil's proof uses $\bar C\times\bar C$, and Connes proposes the square of the arithmetic site as its analogue \cite{ConnesSquare}; these are analogies, not identities (\note{08\_} §§3--4).

For Weil~II the bridge produced a purity criterion, corrected lemmas and a list of misprints. For $\zeta$ it produced a precise form of the weight question (Question~\ref{q:weights}), which for the square is Manin's question of absolute Descartes powers of $\operatorname{Spec}\ZZ$ \cite{Manin}.
My assessment: the programme's even-power computations supply the positivity half of Deligne's mechanism in Weil~II~§1.5, and the other half is the uniform pole bound. Since that bound is equivalent to $\max\re\rho\le\tfrac12$ on finite exponent data, I think the decisive work on this bridge is an independent source for it. The starting point is the programme's own statement (\lbl{TWC9}, \lbl{TWC11}) that its computations leave room for one.

\subsection{Nyman--Beurling in three topologies}

The same kind of family, the cover discrepancies or dilation tests, is studied in three settings: in $\B$, in $L^2(0,\infty)$ with two-sided tests, and in $L^2(0,1)$ and Noor's Hardy space with one-sided tests. In $\B$ the family generates the zero-jet ideal unconditionally, already for the degrees $2^a3^b$ when $t>(\log2)(\log3)/(4\pi)\approx0.0606$ (Proposition~\ref{prop:NB}). In $L^2(0,\infty)$ two-sided tests are dense unconditionally (Wiener's $L^2$ Tauberian theorem \cite{Wiener}). One-sided density in $L^2(0,1)$ is equivalent to RH (Beurling \cite{Beurling}, Báez-Duarte \cite{BaezDuarte}; Burnol's co-Poisson theory \cite{Burnol}), and Noor's Hardy-space form is also equivalent to RH \cite{Noor}. The three settings differ in topology, support and test class, and RH is the statement about the third (theorems; \note{09\_} §4(c), \note{16\_}).
The bridge suggests a comparison map between the settings, with its kernel identified, which would show where the content of RH enters.
My assessment: the contrast is instructive, since the same family is dense unconditionally in one topology and exactly under RH in another. The sparse families of Question~\ref{q:sparse} probe which features of the topology matter.

\subsection{Connes, Meyer and Connes--Consani}

This bridge carries several maps: from Meyer's quotient $Q$ \cite{Meyer} to Connes' weighted Hilbert quotients \cite{Connes}; from the programme's doubled source to Connes--Consani's three-point base $Y$ \cite[§5]{CCP1}; the programme's adelic periodization complex; from Connes--Consani's absolute twistor line \cite{CCtwistor} to the elliptic curve $y^2=x^3-x$; and from the reversed prime knots of \cite{CCknots} to the conjugate character.
These are established:
\begin{itemize}
\item The kernel of the first map is exactly the classes vanishing on all critical-line jets, so it is injective if and only if RH holds; the loss happens in the Hausdorff quotient (Section~\ref{ssec:receivers}; register S26, S28).
\item For sheaves pulled back from the base, the second map induces an isomorphism on cohomology, with supports (Section~\ref{sec:other}; \lbl{DCP7.2}, \note{39\_}). In the trivial finite-unit sector it gives $H^*(\mathbb{P}^1_{\mathbb{F}_1},\Omega)$, whose $H^1$ is Connes--Consani's spectral realisation of the zeros.
\item The adelic periodization complex is formal: after degree-zero coinvariants and in the algebraic representation category, $C\simeq\B_{\mathrm{pr}}[1]\oplus Q[0]$ equivariantly (\lbl{OMS7A.1}; \note{40\_} §4).
\item The twistor line is linked to the CM curve through its $\mathbb{F}_{1^2}$-points (Section~\ref{sec:other}).
\item The prime knots carry the conjugate character when reversed, and for primitive $\chi$ the functional equation puts the anti-character on the other side at the phase $W(\chi)$, with $W(\chi)W(\bar\chi)=1$ (\note{22\_} Proposition 22.8). These are theorems, and their reading as anti-numbers is a proposal.
\end{itemize}
Four further identifications also hold: Connes' harmonic distribution of §VIII, done with the full Gram matrix, is the programme's \lbl{FGR}; spectral synthesis on $Q$ is a local-description theorem in Krasichkov-Ternovskii's sense \cite{Krasichkov}; the positive transfer forms before the quotient are the Bochner--Schwartz theorem in Mellin coordinates \cite{Vladimirov}; and the divisor potential of Section~\ref{ssec:receivers} belongs to the Littlewood-lemma family \cite{BalazardSaiasYor,SekatskiiBeltraminelliMerlini} (identities of objects and theorems, as stated in the programme and audited).

For Connes' approach the bridge suggests an exact account of what the Hausdorff completion loses, and a jointly faithful family of quotients over all lines $\re s=\sigma$ (\lbl{VWR}). For the programme it suggests a dictionary with the established spectral realisations.
My assessment: the comparison of Connes' Hilbert quotients with Meyer's Fréchet quotient is the cleanest structural statement here. The comparison of the classification of positive forms with Meyer's work and with the literature on Weil positivity has still to be made, and I would make it first.

\subsection{Clocks, Bost--Connes, Beurling and monoids}

The programme's normalised clock operators $S_n=\sqrt n\,V_n$ satisfy $S_mS_n=S_{mn}$ and $S_m^*S_n=S_{n/d}S_{m/d}^*$ ($d=\gcd(m,n)$), the relations of the semigroup crossed product $C(\widehat{\ZZ})\rtimes\NN^\times$ (Laca--Raeburn \cite{LacaRaeburn}). So the author's global quotient of all prime clocks is the space on which the Bost--Connes system acts (\lbl{P14}; \note{01\_} §4; theorem). The partition function of that system is $\zeta(\beta)$, with its phase transition at the pole $\beta=1$ \cite{BostConnes}. Its link with the ledger of the pole at $1$ in Proposition~\ref{prop:poles} is an interpretation, for which a map remains to be constructed.
The timed-prime reconstruction is Beurling's generalised-number construction $D=\exp_*(R)$ specialised to the actual primes (\lbl{TP9--TP13}; identity).
At $a=1/q$ the copy's Theorem~E is the factoriality criterion for $\{n\equiv1\bmod q\}$, a Krull monoid with class group $(\ZZ/q)^\times$. The even-sheet test is the same criterion after dividing the class group by $\pm1$, and it coincides with half-factoriality of the one-sided monoid, the Krull-monoid analogue of Carlitz's class-number-two theorem (\cite{BaginskiChapman,Carlitz}; Section~\ref{sec:sheets}; theorem).

For the theory of monoids the bridge yields a freeness test for every monoid of integers (Theorem~\ref{thm:formallog}). For generalised primes it shows that regularity of the counting function alone does not force the analogue of RH (item N-Beurling).
My assessment: the monoid results are of independent interest as elementary number theory, Theorem~\ref{thm:formallog} and Corollary~\ref{cor:firstneg} most clearly. The novelty of the theorem was checked by one search, and that of the corollary was not searched.

\subsection{Two-prime density and mean-periodic functions}

\lbl{MCL3.4}, \lbl{IH6} and \lbl{PL3} rest on one fact: $(\log p)\ZZ+(\log q)\ZZ$ is dense in $\R$. This is the multiplicative form of the two-period phenomenon of mean-periodic functions \cite{SchwartzMP}. The density holds because $\log p/\log q$ is irrational, and the link with mean-periodic functions is a comparison of mechanisms.
The bridge explains why two primes determine every zero while one prime leaves the aliases $\rho+2\pi ik/\log p$ (Section~\ref{sec:other}).
My assessment: for mean-periodic functions spectral synthesis is a theorem: every closed translation-invariant subspace of $C(\R)$ is the closed span of the exponential monomials $x^ke^{\lambda x}$ it contains \cite{SchwartzMP}. Under $u=e^x$ these monomials become the functions $u^{\lambda}(\log u)^k$ that pair with the jets at the zeros, so I think the transfer of Schwartz's theorem to the programme's jet calculus is a direct next step, and one whose outcome can be checked.

\subsection{The anomaly dictionary (a proposal)}

The author's physical picture is matched term by term with exact objects (\note{22\_} §1). The objects and the facts proved about them hold without assuming RH; the matching itself is a proposal.
\begin{center}\small
\begin{tabularx}{\linewidth}{@{}l X@{}}
\toprule
Author's term & Exact object\\
\midrule
anti-number & the clock run backwards with conjugation, $g\mapsto\cl{g(1/u)}/u$ (Mellin shadow $s\mapsto1-\bar s$)\\
chirality & the mirror $A(s)=-\bar s$ through the pole at $0$\\
the doublet & the orbit $\{\rho,\rho^\#,\rho^\#-1,\rho-1\}$ of an off-line zero in the two-line system\\
$\ZZ_2$ anomaly & the two lines as the two sheets of one double cover, deck transformation the central $-1$ of $SU(2)$\\
$\ZZ_1$ anomaly & the root-number phase $W(\chi)$, cancelled by the anti-sector: $W(\chi)W(\bar\chi)=1$\\
inflow & the bulk between the lines, which holds exactly the off-line zeros with $\re\rho<\tfrac12$ and their mirrors\\
the support $\tau$ & the pivot of the swing between $0$ and $1$: on the zeta side the reflection $s\mapsto1-\bar s$ about $\tfrac12$\\
\bottomrule
\end{tabularx}
\end{center}
In these terms RH has three exact forms. Every zero is its own anti-number. The bulk is empty. And the Weil functional is nonnegative on every annihilation $g*g^\#$ (Weil's criterion); the value of the annihilation at the unit, $\int|g|^2du$, is positive with or without RH.
My assessment: the dictionary is useful because every entry is an exact object. Its test is whether a statement first seen in the physical language yields a new statement about these objects.

\subsection{The \texorpdfstring{$\mathbb{F}_1$}{F1} context}

The author describes $\tau$ as carrying $\ZZ_1$ and no $\ZZ_2$. This is realised at Connes--Consani's generic point $\eta$, which every homeomorphism fixes and where no exchanged label can sit, although the invariant local data there still contain the sign $\varepsilon$. The absence of $\ZZ_2$ thus has two exact readings: no exchanged label, which holds at $\eta$; or no sign at all, $\varepsilon=1$, which holds among field-valued points exactly in characteristic $2$, through the retraction $\mathbb{F}_{1^2}\to\mathbb{F}_1$.
The labels $\ZZ_n$ follow the principle that an $\mathbb{F}_{1^n}$-structure is a free $\ZZ/n$-action (Kapranov--Smirnov, via Connes--Consani--Marcolli \cite{CCMfun}, §3; a reading). The points of $\mathbb{P}^1$ over $\mathbb{F}_{1^2}$ are $\{0,\infty,\pm1\}$, and the double cover branched at them is $y^2=x^3-x$.
The Eulerian sheets $a=1/q$ of the even lattice family are exactly those with $(\ZZ/q)^\times=\{\pm1\}$, the units of $\mathbb{F}_{1^2}$, and exactly those on which the one-sided monoid is half-factorial (Proposition~\ref{prop:evensheets}); reading this as $\mathbb{F}_1\to\mathbb{F}_{1^2}$ is an interpretation.
The global quotient of all prime clocks is $\widehat{\ZZ}$ (Proposition~\ref{prop:clocks}), the Bost--Connes space, which Connes and Consani treat as an $\mathbb{F}_1$ object (a reading). The split-zero semiring $G(\ZZ)$ adds a Boolean point below the generic point of $\operatorname{Spec}\ZZ$, where the sign is forgotten, while Connes and Consani add a closed point $\infty$ glued along the generic point (a reading).
The question at which the programme's analysis of weights arrived can be stated without its vocabulary. Does the reconstructed arithmetic carry either (a) a square with a Künneth map $H^1\otimes H^1\to H^2$ and an auxiliary object with discrete weights, or (b) an independent cohomological source for a pole bound uniform in the tensor power (Section~\ref{sec:deligne})? For the square this is Manin's question of absolute Descartes powers of $\operatorname{Spec}\ZZ$ \cite{Manin}.
My assessment: the reading with mathematical content of its own is the coincidence of the Eulerian even sheets with the units of $\mathbb{F}_{1^2}$, which deserves a statement independent of the reading.

\subsection{The author's other workbenches}\label{ssec:overlaps}

The author maintains workbenches on the Erdős--Straus conjecture (ES), on Yang--Mills (YM, which also holds the Navier--Stokes (NS), $S^6$ and Jacobian-map material), on Erdős problem 817, and on the Collatz problem. Companion records in the same folder treat these workbenches and the other sides of the edges below. The table lists a selection of the edges located from the zeta side, audited in the notes \note{21\_}, \note{28\_} and \note{34\_} and, where stated, in the companion records. The Collatz side of the last edge was located in an external record, the Clankers' Zenodo record 19900461; the author's Collatz workbench itself was not read for this record, and the companion record treats it. The status column says what the audits established: \emph{verified} means read and checked, with the note or record cited; \emph{located} means found and described, not checked.

\begin{center}\small
\begin{longtable}{@{}p{0.2\linewidth}p{0.47\linewidth}p{0.27\linewidth}@{}}
\toprule
Edge & Shared object or lemma & Status\\
\midrule
\endhead
NS $\to$ ES & the torus endomorphism $J=\begin{psmallmatrix}3&1\\1&5\end{psmallmatrix}$ of the NS manuscript, used as an operator in 21 ES statements; exact aliasing on finite character groups; sharp constant $1/\sqrt{4+2\sqrt2}$ for the $\sqrt2$ directions & verified (register S40); joint averages need the complete preimage (N-X)\\
NS $\to$ YM & NS velocity as a reducible $SU(2)$ connection; $-2\sum\operatorname{tr}\mathcal{F}_{ij}^2=\lambda^2|\operatorname{curl}u|^2$ & verified, conditional on imported NS bounds and on YM sections not read; zero-quotient sequence only as the coupling tends to $0$\\
zeta $\to$ YM & the Gram sandwich from a relative column error & verified (register S41)\\
Jacobian map $\to$ NS kinematics $\to$ YM & $U=DF^{-1}(V\circ F)$, the escaping curve, the material tensor; YM Theorem~14.2 & verified: Lemma~\ref{lem:keller}; Theorem~14.2 correct conditional on fixed-box weak-coupling limits (N-X)\\
$S^6\to$ YM & the imaginary period block $B$ of the $S^6$ period matrix enters only through $D=\det B$ & verified (register S53)\\
$S^6\to$ ES & the integral monodromies of the $(3,4,\infty)$ period system give decorated affine-torsion covers of the ES divisor atlas & located in version~2; the orbits of the covers are now verified at every level in the Erdős--Straus record\\
ES $\leftrightarrow$ zeta & the split-zero semiring $G(R)=R\sqcup\{\tau\}$ with its unique Boolean character & verified (the zeta proof checked by hand)\\
zeta $\to$ ES & ES family-C results use zeta \lbl{NG20}--\lbl{NG21} as theorem dependencies & located (a dependency, not re-proved in ES)\\
ES $\leftrightarrow$ zeta & a four-dimensional Keller map $P$ with $\det DP=-2$, transported into the split-zero weighted conductor; the shared file is byte-identical in both repositories & $\det DP=-2$ verified; the rest located\\
NS $\to$ zeta & the released NS field as arithmetic input (Mellin transports, residue at $s=-1$) & located; depends on the imported NS existence theorem, whose independent verification the zeta side records as unfinished\\
817 $\to$ zeta & finite word sources of problem 817 mapped into the split-zero supported quotient (typed maps written; no zeta statement is entered) & located in version~2; one of the ten interface notes is now verified in the companion record\\
Collatz $\leftrightarrow$ ES $\leftrightarrow$ zeta & the support-idempotent homogenisation $(x,h)\mapsto(ax+bh,h)$ & located in version~2 (same formula on all three sides); verified in part in the companion record\\
\bottomrule
\end{longtable}
\end{center}

One Collatz item is elementary and can be checked here. With $T(n)=4n+1$, $\theta(n)=3n+1$ and the odd step $U(n)=(3n+1)/2^{v_2(3n+1)}$, one has $\theta(T(n))=12n+4=4\theta(n)$, hence $U(T^jn)=U(n)$ for all $j$ (ES results bench, R05).
The Erdős--Straus workbench states that it does not claim a proof or counterexample of the conjecture, and the Collatz workbench that it is not a claim that the Collatz conjecture has been solved; the Yang--Mills and Navier--Stokes material states its own scope limits, recorded in N-X.

My assessment: the shared linear algebra (Gram sandwiches, Schur complements) is the most reusable crossing, since it is exact and needs no identification of objects. The edges that carry a statement about $\zeta$ into another workbench, or back, still need typed maps with verified content.
